\scanpage{36}
(évident sur (8)), qui \deleted{\unsure{se définit}}\inserted{peut se décrire} comme l'ens des $(p,\sigma)$ $(p\in\hat{\mathbb Z}^*,\ \sigma\in\unsure{\mathfrak S_3|\mathcal S_3})$ tels que l'on ait
\[
(13)\qquad l_{\sigma(0)}^p\ l_{\sigma(1)}^p\ l_{\sigma(\infty)}^p=1
\]
Quand on se borne à $S\mathcal E$ au lieu de $S\widehat{\widehat{\mathcal E}}$, de sorte qu'on doit avoir $p\in\mathbb Z^*=\{\pm1\}$, \deleted{on trouve que les} i.e. l'ens des $(p,\sigma)$ (\unsure{abstrayant|satisfaisant} la condition (13)) n'est autre que $\deleted{\unsure{\mathbb Z}}\inserted{\{\pm1\}}\times\unsure{\mathfrak S_3|\mathcal S_3}\simeq\unsure{\mathfrak S_3|\mathcal S_3}\times\mathbb Z_2\simeq\Gamma$, on voit que $(p,\sigma)$ satisfait (13) sss $(\sigma,p)\in\Gamma_P$. Ainsi, $\tilde\Gamma_P\subset S\mathcal E\subset S\widehat{\widehat{\mathcal E}}$ est caractérisé comme le sous-groupe \inserted{de $S\mathcal E$} des autom. de $\pi$ \deleted{\unsure{ou} $\unsure{\widehat{\widehat\pi}|\hat\pi}$} qui \unsure{invarient} $L_0,L_1,L_\infty$, affectés des \deleted{multiplicateurs}\inserted{``rectificateurs''} $g_0=1,g_1=1,g_\infty=1$.

Sans doute il est vrai que le $\unsure{=|\bar m}$ résultat est valable dans $S\widehat{\widehat{\mathcal E}}$, i.e. que la relation (13) implique forcément $p\in\{\pm1\}$. Il revient au $=$ de dire \inserted{(quitte à conjuguer par un $\tilde\sigma_i$ --- \unsure{ou} $\sigma\in\Gamma_P$)} que la relation
\[
(14)\qquad l_\infty^p\ l_1^p\ l_0^p=1
\]
(pour un $p\in\hat{\mathbb Z}^*$ fixé) implique $p=1$.

Exercice!
\mspage{35}
